\subsection{Comparison with earlier bounds}
\label{subsec:literaturecomparison}

The union, multiplicity-two mass, and multiplicity-two distinct-location
conclusions improve identified unconditional bounds for the same counting
statistics. The union bound improves Lamzouri's
Theorem~1.1 \cite{Lamzouri}; the two multiplicity-two bounds improve
the displayed formulas in Waddle's public research draft \cite{Waddle}.
The respective improvements, as differences of proportions, exceed
\[
 0.000934060819863356,\qquad
 0.001118156681831034,\qquad
 0.000469272388685036.
\]
Multiplying these differences by $100$ gives percentage-point gains.
The comparisons use the exact expressions, not subtraction of rounded
headline percentages.

To state the comparisons compactly, put
\begin{equation}\label{eq:literatureconstants}
 C_0=2-R_0,\qquad C_1=\frac{3-R_0}{2},\qquad
 B=(3-2\sqrt2)(R_0-1),\qquad C_2=1-2B.
\end{equation}
Thus $C_0=0.672500703679411\ldots$ is the Montgomery--Taylor
simple-critical constant and $C_1=0.836250351839705\ldots$ is the
corresponding distinct-zero constant. Alp\"oge and Furman
\cite[Theorem A]{AF} established these bounds unconditionally by an
inertia and rank--trace argument for a finite compression of Weil's
Hermitian form. Lamzouri \cite{Lamzouri} recovered them through a
Hilbert-space argument and obtained the union constant
\[
 C_2=\frac{1+2\sqrt2+2C_0}{3+2\sqrt2}
       =0.887620008173354\ldots.
\]

\begin{table}[ht]
\centering
\small
\setlength{\tabcolsep}{4pt}
\renewcommand{\arraystretch}{1.6}
\begin{tabular}{@{}p{.10\textwidth}p{.43\textwidth}p{.40\textwidth}@{}}
\hline
Ratio & Present lower bound & Named comparison\\
\hline
$U/N$
& $1-2E$\par $>0.888554068993217696$
& Lamzouri \cite{Lamzouri}:\par
  $C_2=0.887620008173354\ldots$\\
$M_{\le2}/N$
& $\dfrac{3-T_1}{2-a_1}$\par $>0.837368508521536857$
& Waddle \cite{Waddle}:\par
  $C_1=0.836250351839705\ldots$\\
$D_{\le2}/N^{\rm d}$
& $1-\dfrac{\beta}{8-2T_1-a_1}$\par $>0.939197203148497742$
& Waddle \cite{Waddle}:\par
  $1-\dfrac{R_0-1}{2(4-R_0)}$\par
  $=0.938727930759812\ldots$\\
$n/N$
& $\dfrac{2-R_0-d_0}{1-a_0}$\par $>0.673058325315610967$
& Devine \cite{Devine}:\par
  $673399/10^6=0.673399$\par
  (stated unconditional theorem)\\
$N^{\rm d}/N$
& $\dfrac{3-R_0-d_0-a_0}{2(1-a_0)}$\par $>0.836529162657805483$
& Alp\"oge--Furman and Lamzouri\par
  \cite{AF,Lamzouri}:\par
  $C_1=0.836250351839705\ldots$\\
\hline
\end{tabular}
\caption{Five separate global lower limiting proportions. Each exact
expression in the middle column bounds the corresponding lower limit
from below; the strict inequalities compare these expressions with
shorter decimals. All external statements listed here are unconditional.
Waddle's source is an unrefereed public draft and Devine's is a public
preprint. The last row is a comparison with the stated foundational
bounds, not a claim of an overall distinct-zero record.}
\label{tab:literaturecomparison}
\end{table}

Here $n$ counts simple critical zeros, whereas $U$ counts with
multiplicity the union of simplicity and critical-line location.
The denominators $N$ and $N^{\rm d}$ are different: $M_{\le r}/N$
is a fraction of zero mass, and $D_{\le r}/N^{\rm d}$ is a fraction
of distinct complex zero locations. Locations with the same imaginary
part are not identified. These conventions also match the external
statements in Table~\ref{tab:literaturecomparison}. A dyadic statement
for additive counts gives the corresponding cumulative lower limit by
summing sufficiently high dyadic pieces of $(0,X]$; the bounded initial
remainder disappears after division by the same denominator.

For every fixed integer $r\ge3$, Waddle's displayed mass and
distinct-location bounds are respectively
\[
 1-\frac{(r+1)B}{r-1},\qquad
 1-\frac{B}{r-1-rB}.
\]
Our corresponding formulas replace $B$ by $E$, and the proved
inequality $0<E<B<1/2$ makes both improvements strict. The mass
gain is $(r+1)(B-E)/(r-1)$; the location comparison follows from
the strict increase of $x/(r-1-rx)$ on the relevant interval.
This conclusion uses the retained multiplicity budget. It is not a
formal consequence of a union bound: a multiset supported on real
locations of multiplicity $r+1$ has full union mass and no mass at
multiplicity at most $r$.

The fixed-cutoff statement is an improvement over a specified family.
It does not assert the best available bound for every cutoff.
Simoni\v{c}, Trudgian and Turnage-Butterbaugh \cite[Theorems 4 and 7]{STT}
give unconditional exponential estimates for the number of locations
of exact multiplicity $j$. Weighting by $j$ and summing the tail gives
mass bounds tending to one as the cutoff grows, whereas the displayed
mass family here tends to $1-E<1$. Consequently the exponential-tail
bound is stronger for sufficiently large fixed cutoffs.

The simple-critical conclusion refines Ainta's public bound
\cite{Ainta}, whose exact expression is
\[
 \frac{1345000C_0-2680}{1340003}
       =0.673008527927779\ldots.
\]
It remains below Shi's stated research-draft candidate
$0.6733169771424713\ldots$ \cite{Shi} and Devine's stated
$0.673399$ \cite{Devine}. The latter preprint's numerical proof has
not been independently reconstructed here. These are external
comparison targets, not premises of our proof, and we claim no overall
simple-critical or distinct-zero record. Zeta Lab's public report
\cite{ZetaLab} also records an eight-point value
$0.67305298298962888\ldots$, distinguishing an interval-certified
local inequality from its Lean implication conditional on that
inequality. The public union extension \cite{Union} supplies another
explicit treatment of the distinction between union and intersection.

Stronger numbers with additional hypotheses must be compared separately.
The union bounds $8/9$ in \cite[Remark 2]{BGSTI} and
$(2+C_0)/3=0.890833567893\ldots$ discussed in \cite[Remark 1.2]{Lamzouri}
use a narrow-box assumption on the horizontal locations of all
sufficiently high dyadic zeros; they are not unconditional union
comparators. Under RH, Gon\c{c}alves, de Laat and Leijenhorst
\cite[Theorem 1]{GdLL} prove
$\liminf M_{\le2}/N\ge0.9614$. Their cutoff-three value $0.9787$
also carries the theorem's additional qualification concerning Maple
summation of specified series. RH-dependent pair-correlation estimates
in \cite{CGdL} include $0.6792$ for simple zeros and $0.8477$ for
distinct zeros. Devine's claims beyond $0.6792$ likewise require
additional density or horizontal-profile assumptions \cite{Devine}.
None of these numbers is an unconditional obstruction to the bounds
proved here.


\subsection{Relation to earlier methods}
\label{subsec:literaturenovelty}

The analytic framework and several finite ingredients have established
antecedents. Montgomery's pair-correlation theorem \cite{Mon73} and
the Montgomery--Taylor window reported in \cite{Mon75} underlie the
constant $R_0$; the associated extremal problem is treated in
\cite{CCLM}. Aryan \cite{Aryan} obtained an unconditional form of
the relevant Fej\'er-kernel calculation, and Baluyot, Goldston,
Suriajaya and Turnage-Butterbaugh \cite{BGST} proved the published
full-height pair-correlation formula used here. Lamzouri's
finite-multiset formulation and removal of the rational pair weight
\cite{Lamzouri} provide the direct analytic interface. Wang
\cite{Wang} supplies the separate short-height input, with support
strictly smaller than the fixed interval exponent. We prove no new
pair-correlation theorem and do not insert a height-dependent test
into a fixed-test error estimate.

Ainta \cite{Ainta} introduced the clipped spectral stability
correction and quantitative overlapping seven-point windows used in
this line of refinements. The Fourier perturbations of the cosine
density and nonuniform pair weights with distance-capacity constraints
in the inspected Ojha and Shi drafts \cite{Ojha,Shi} are also prior
work. Shi's threshold-one trace--energy envelope, pressure rescaling,
and two-certificate supporting-plane deduction are directly relevant
antecedents. In particular, neither weighted overlap nor removal of
the elementary threshold-one cap is claimed as a new idea here.
The arbitrary-parameter scalar multiplicity framework and its
uncorrected low-multiplicity consequences are credited to the
two-trace development and Waddle's displayed deductions
\cite{AF,Waddle}.

The finite contribution here is the use of one common averaged witness
to transfer a whole-domain local Gram inequality directly to the
clipped spectral defect while preserving its pressure floor. The
uniform coefficient estimate depends on the clipping threshold.
For the seven-point weights used here, its admissible cap is
$40t^2/2401$. This dependence is material: the new local floor
$a_1=7/400$ is unavailable through this cap at $t=1$, but is
admissible at $t=2$ and at the union threshold
$t=1+\sqrt{2(1-a_1)}$. Thus the stronger local certificate is matched
to the count being proved. Retaining the resulting simple-Gram
correction in the separate per-location multiplicity costs then
produces the mass and distinct-location bounds from that same
certificate. The elementary variational identity for clipping is
part of the proof; the gain comes from its common-witness application
and the uniform capacity estimate.

The required two-mode profile, its normalization and energy cost,
and the whole-domain seven-point inequality are established explicitly.
The local interval engine is an attributed adaptation of Ainta's
MIT-licensed implementation, using Arb ball arithmetic \cite{Johansson}.
The separate cosine row for the simple-critical and total-distinct
conclusions is Ainta's completed input. Our reproduced executions
and exact scalar calculations certify the stated finite numerical
premises; they do not replace the finite operator argument or the
passage to zero limits. 

Finally, bounded-potential versions of the common-witness argument
give additional finite inequalities, but the general telescoping
subaction idea already occurs in the public literature, including
\cite{Devine}. Partial larger-window certificates are identified by
their domains. They supply no global zero-count conclusion unless
the entire required configuration space has been covered.
